\chapter{กรณีศึกษาเชิงลึก: การสร้าง Mr. Deep\_xAI และ J.A.R.V.I.S. Protocol}

การสร้างบอทปัญญาประดิษฐ์ที่มีปฏิสัมพันธ์กับผู้ใช้แบบเรียลไทม์ ต้องผสานทั้งความสามารถในการมองเห็น (Vision), การพูดคุยสื่อสาร (Speech Synthesis), และการวิเคราะห์ตรรกะเชิงวิทยาศาสตร์

\section{การสกัดภาพและการตรวจสอบฉาก 3 มิติแบบเรียลไทม์}

บอท Mr. Deep\_xAI ตรวจสอบสภาพแวดล้อม Three.js โดยตรงผ่านการดึงเฟรมภาพจาก \texttt{canvas.toDataURL("image/jpeg")} นำเข้าสู่มัลติโมดัลวิทัศน์ เพื่อตรวจสอบตำแหน่งลูกตุ้ม การต่อขั้วสายไฟ หรือค่าที่แสดงบนหน้าปัดมิเตอร์

\begin{teachingplan}{การตรวจจับข้อผิดพลาดในวงจรไฟฟ้า 3 มิติ}
เมื่อผู้ใช้เชื่อมต่อสายไฟผิดขั้ว บอทจะตรวจพบว่าค่าศักย์ไฟฟ้าที่จุดทดสอบไม่สอดคล้องกับสมการเคอร์ชอฟฟ์ (Kirchhoff's Laws) และแจ้งเตือนผ่านเสียงพูดทันที
\end{teachingplan}

\section{ระบบเสียงสองทางสำเนียงสุโขทัยแท้}

เพื่อสร้างความรู้สึกเป็นกันเองและมีชีวิตชีวา บอทใช้เอนจิน Google Neural TTS ร่วมกับโมเดลภาษาที่ได้รับการปรับแต่งสำเนียงไทยสุโขทัย โดยมีหางเสียงและการลงท้ายประโยคที่เป็นเอกลักษณ์เฉพาะถิ่น

\section{วงจรตัวกรองเสียงสื่อสารอวกาศด้วย Web Audio API}

สัญญาณเสียงของบอทได้รับการประมวลผลผ่านตัวกรอง DSP Filter Chain ในเบราว์เซอร์ เพื่อให้มีมิติเสียงเสมือนหุ่นยนต์คู่หูระดับไฮเทค:

\begin{itemize}
    \item \textbf{BiquadFilter (Bandpass 450--3500 Hz):} ตัดความถี่ส่วนเกินเพื่อให้ได้เสียงวิทยุสื่อสาร
    \item \textbf{ConvolverNode (Impulse Response):} จำลองความก้องของชุดเกราะห้องทดลอง
    \item \textbf{StereoPannerNode:} ปรับตำแหน่งมิติเสียงตามการเคลื่อนไหวของอวาตาร์บนหน้าจอ
\end{itemize}
